% SPDX-FileCopyrightText: © 2026 Idiap Research Institute <contact@idiap.ch>
%
% SPDX-FileContributor: Andre Freitas <andre.freitas@idiap.ch>
% Neuro-symbolic AI Group
%
% SPDX-FileContributor: Danilo Gusicuma <danilo.gusicuma@idiap.ch>
% Neuro-symbolic AI Group
%
% SPDX-License-Identifier: MIT
\section{Formal Complexity Derivation}

\label{sec:formal-derivation}

This section is generated from structured definitions, assumptions, propositions, calculations, and unresolved obligations. It states a local paper-level object and does not claim the global target theorem.

\subsection{Definitions}

\paragraph{H\_N.} A source carrier language, truth-table family, meta-complexity predicate, or proof-mediated object indexed by source length N.

\[
H_N\subseteq\{0,1\}^{N}\quad\text{or}\quad H_N:\{0,1\}^{N}\to\{0,1\}
\]

\emph{Valid when:} The carrier is explicitly constructed or specified without nonuniform hard-row advice.

\paragraph{SAT\_m.} The target SAT language or SAT truth-table slice at input length m, used as the reference NP-complete target.

\[
\mathrm{SAT}_m=\mathrm{SAT}\cap\{0,1\}^m
\]

\emph{Valid when:} The reduction or transfer states how source length N maps to target length m.

\paragraph{a,q,beta.} The source lower-bound exponent, overhead exponent, and target-query-length exponent in a transfer inequality.

\[
\mathrm{SIZE}(H_N)\le N^{q+\beta c+o(1)}\quad\text{from}\quad \mathrm{SAT}\in\mathrm{SIZE}(n^c)
\]

\emph{Valid when:} All padding, decoding, verifier, and advice costs are counted in the same circuit model.

\subsection{Assumptions}

\begin{itemize}
\item The source carrier H\_N and target class are named before any transfer claim is made.
\item The circuit model is fixed throughout the transfer, including uniformity and advice conventions.
\item The transfer from SAT circuits to source-carrier circuits has an explicit exponent envelope q+beta*c+o(1).
\item Barrier status is route-specific: naturalness, relativization, algebrization, and advice leakage are attached to the actual theorem shape.
\end{itemize}

\subsection{Local Propositions}

\paragraph{Bounded-window transfer obstruction.} A barely-superlinear source lower bound excludes only the target exponents that fit below the transferred upper-bound envelope.

\[
\mathrm{SIZE}(H_N)>N^a\ \wedge\ \mathrm{SAT}\in\mathrm{SIZE}(n^c)\Rightarrow \mathrm{SIZE}(H_N)\le N^{q+\beta c+o(1)}\quad\Longrightarrow\quad a>q+\beta c
\]

\emph{Proof steps.}

\begin{itemize}
\item Assume SAT has nonuniform circuits of size n\textasciicircum{}c in the stated model.
\item Apply the route transfer to build source-carrier circuits of size N\textasciicircum{}\{q+beta c+o(1)\}.
\item Contradict the source lower bound only for those c with a>q+beta c.
\item If beta>0, a fixed a excludes at most a bounded exponent window.
\end{itemize}

\emph{Valid when:} The source and target circuit models, length map, and overhead exponents are unchanged across the implication.

\paragraph{All-exponent lift gate.} A route can support NP not subset P/poly only if it preserves the all-exponent quantifier or supplies a cofinal indexed family.

\[
\forall c\ \exists N_0\ \forall N\ge N_0:\quad a_c>q_c+\beta_c c
\]

\emph{Proof steps.}

\begin{itemize}
\item P/poly exclusion requires ruling out every fixed polynomial exponent c for SAT or an NP-complete target.
\item For each c, the source lower bound must outrun the transferred upper bound.
\item If the same source theorem has beta>0 and fixed a, large c survives.
\item Thus promotion needs beta=0 with positive margin or a cofinal parameter family with no advice leakage.
\end{itemize}

\emph{Valid when:} The parameter family is uniform or advice-accounted and valid on target-relevant lengths.

\paragraph{Searchable certificate gate.} A meta-complexity or proof-object bridge is useful only when certificates are frequent, searchable, sound, and advice-free without becoming a large constructive natural property.

\[
\exists u,\pi\ \bigl(|\pi|=O(m)\wedge V(u,\pi)=1\bigr)\Rightarrow CC(f_u)>2^{\epsilon m}
\]

\emph{Proof steps.}

\begin{itemize}
\item State the verifier, decoder, certificate length, and search schedule.
\item Prove that accepted pairs exist in every required gap interval.
\item Prove soundness from certificate acceptance to the claimed circuit lower bound.
\item Audit whether the accepted set is large and constructive, or whether row selection leaks advice.
\end{itemize}

\emph{Valid when:} The verifier and search live inside the intended target class and do not choose hard rows nonuniformly.

\subsection{Worked Calculation Interface}

\paragraph{Direct-cluster soundness falsifier.} Odd-charge Tseitin-style parity ambiguity is a finite falsifier for local consistency claims that do not bind the target row semantically.

\[
\mathrm{LocalConsistent}(C)\not\Rightarrow \mathrm{DirectClusterSound}(C,k)
\]

\emph{Valid when:} The tested proof system exposes only local projections or bounded-width proof views and lacks a global parity-binding axiom.

\subsection{Unresolved Obligations}

\begin{itemize}
\item Instantiate one named proof system or meta-complexity carrier and compute q, beta, advice cost, and certificate-search exponent.
\item Prove the direct-cluster or parity-binding lemma, or exhibit the odd-charge Tseitin counterexample for that named system.
\item Run a prior-art boundary check against hardness magnification, MCSP/MKtP, natural-proofs, and proof-complexity transfer literature.
\item State whether the surviving result is a positive bridge theorem, a scoped no-go theorem, or a demoted route diagnostic.
\end{itemize}

\subsection{Non-Claims}

\begin{itemize}
\item The paper does not prove P != NP.
\item The paper does not prove SAT is outside P/poly.
\item A route-level transfer obstruction is not a target separation theorem.
\item Artifact evidence documents the run state; external novelty still requires a prior-art boundary.
\end{itemize}
